\section{FLAT MODULES AND ``TOR'' FUNCTORS}

For the benefit of readers conversant with Homological Algebra (*), we shall indicate quickly how the theory of flat modules is related to that of Tor functors.

PROPOSITION 1. Let E be a right A-module. The following four properties are equivalent:

\begin{enumerate}
\def\labelenumi{(\alph{enumi})}
\item
  E is flat.
\item
  For every left A-module F and every integer \(n \geqslant 1\) , \(\operatorname{Tor}_n^{\mathbf{A}}(\mathbf{E}, \mathbf{F}) = 0\) .
\item
  For every left A-module F, \(\operatorname{Tor}_{1}^{\mathbf{A}}(\mathbf{E},\mathbf{F}) = 0\) .
\item
  For every finitely generated left ideal a of A,
\end{enumerate}

\[
\operatorname{Tor} _ {1} ^ {\mathbf {A}} (\mathrm{E}, \mathrm{A} _ {s} / \mathfrak {a}) = 0.
\]

We show that (a) implies (b). Let

\[
\dots \rightarrow \mathrm{L} _ {n} \rightarrow \mathrm{L} _ {n - 1} \rightarrow \dots \rightarrow \mathrm{L} _ {0} \rightarrow \mathrm{F} \rightarrow 0
\]

be a free resolution of F. As E is flat, the sequence

\[
\dots \rightarrow \mathrm{E} \otimes \mathrm{L} _ {n} \rightarrow \mathrm{E} \otimes \mathrm{L} _ {n - 1} \rightarrow \dots \rightarrow \mathrm{E} \otimes \mathrm{L} _ {0} \rightarrow \mathrm{E} \otimes \mathrm{F} \rightarrow 0\tag{1}
\]

is exact. As the \(\operatorname{Tor}_n^{\mathbf{A}}(\mathbf{E},\mathbf{F})\) are isomorphic to homology groups of the complex (1), they are zero for \(n\geqslant 1\) . It is trivial that (b) implies (c) and (c) implies (d). We show finally that (d) implies (a). The exact sequence

\[
0 \rightarrow a \rightarrow A _ {s} \rightarrow A _ {s} / a \rightarrow 0
\]

gives the exact sequence

\[
\operatorname{Tor} _ {1} ^ {\mathbf {A}} (\mathrm{E}, \mathrm{A} _ {s} / \mathfrak {a}) \rightarrow \mathrm{E} \otimes_ {\mathbf {A}} \mathfrak {a} \rightarrow \mathrm{E} \otimes_ {\mathbf {A}} \mathrm{A}.
\]

As (d) holds, the canonical homomorphism

\[
\mathbf {E} \otimes_ {\mathbf {A}} \mathfrak {a} \rightarrow \mathbf {E} \otimes_ {\mathbf {A}} \mathbf {A} = \mathbf {E}
\]

is injective, which means that E is flat (§ 2, no. 3, Proposition 1).

Proposition 1 provides a characterization of flat modules which is often useful in applications. We shall restrict ourselves, by way of an example, to giving a new proof of Proposition 5 of § 2, no. 5. If E' and E'' are flat, the exact sequence

\[
\operatorname{Tor} _ {1} ^ {\mathbf {A}} \left(\mathrm{E} ^ {\prime}, \mathrm{F}\right)\rightarrow \operatorname{Tor} _ {1} ^ {\mathbf {A}} (\mathrm{E}, \mathrm{F}) \rightarrow \operatorname{Tor} _ {1} ^ {\mathbf {A}} \left(\mathrm{E} ^ {\prime \prime}, \mathrm{F}\right)
\]

shows that \(\operatorname{Tor}_1^{\mathbf{A}}(\mathbf{E},\mathbf{F}) = 0\) for every left A-module \(\mathbf{F}\) , hence \(\mathbf{E}\) is flat. If \(\mathbf{E}\) and \(\mathbf{E}''\) are flat, the exact sequence

\[
\operatorname{Tor} _ {2} ^ {\mathbf {A}} \left(\mathrm{E} ^ {\prime \prime}, \mathrm{F}\right)\rightarrow \operatorname{Tor} _ {1} ^ {\mathbf {A}} \left(\mathrm{E} ^ {\prime}, \mathrm{F}\right)\rightarrow \operatorname{Tor} _ {1} ^ {\mathbf {A}} (\mathrm{E}, \mathrm{F})
\]

shows that \(\operatorname{Tor}_1^{\mathbf{A}}(\mathbf{E}',\mathbf{F}) = 0\) , hence \(\mathbf{E}'\) is flat.

PROPOSITION 2. Let R, S be two rings, \(\rho: \mathbf{R} \to \mathbf{S}\) a homomorphism and F a left R-module. The following two properties are equivalent:

\begin{enumerate}
\def\labelenumi{(\alph{enumi})}
\item
  \(\operatorname{Tor}_1^{\mathbf{R}}(\rho_*(\mathbf{E}),\mathbf{F}) = 0\) for every right S-module E.
\item
  The left S-module \(\rho^{*}(\mathbf{F}) = \mathbf{F}_{(\mathbf{S})} = \mathbf{S} \otimes_{\mathbb{R}} \mathbf{F}\) is flat and \(\operatorname{Tor}_{1}^{\mathbb{R}}(\rho_{*}(\mathbf{S}_{d}), \mathbf{F}) = 0\) .
\end{enumerate}

Suppose that (a) holds. Taking \(E = S_{d}\) , we see that \(\operatorname{Tor}_{1}^{\mathbb{R}}(\rho_{*}(S_{d}), \mathbf{F}) = 0\) . We show also that \(F_{(S)}\) is a flat S-module. For that, we note that, if E is a right S-module, the additive group \(E \otimes_{S} F_{(S)}\) is identified with \(\rho_{*}(E) \otimes_{R} F\) . Then if there is an exact sequence of right S-modules

\[
0 \rightarrow \mathrm{E} ^ {\prime} \rightarrow \mathrm{E} \rightarrow \mathrm{E} ^ {\prime \prime} \rightarrow 0
\]

we obtain, using (a), an exact sequence

\[
0 \rightarrow \rho_ {*} (\mathrm{E} ^ {\prime}) \otimes_ {\mathrm{R}} \mathrm{F} \rightarrow \rho_ {*} (\mathrm{E}) \otimes_ {\mathrm{R}} \mathrm{F} \rightarrow \rho_ {*} (\mathrm{E} ^ {\prime \prime}) \otimes_ {\mathrm{R}} \mathrm{F} \rightarrow 0
\]

or also

\[
0 \rightarrow \mathrm{E} ^ {\prime} \otimes_ {\mathbf {S}} \mathrm{F} _ {(\mathbf {S})} \rightarrow \mathrm{E} \otimes_ {\mathbf {S}} \mathrm{F} _ {(\mathbf {S})} \rightarrow \mathrm{E} ^ {\prime \prime} \otimes_ {\mathbf {S}} \mathrm{F} _ {(\mathbf {S})} \rightarrow 0
\]

which proves that \(\mathbf{F}_{(\mathbf{S})}\) is flat.

Conversely, if (b) holds, we have first of all, for every free right S-module \(L = S_{d}^{(1)}\) , \(\operatorname{Tor}_{1}^{\mathrm{R}}(\rho_{*}(\mathrm{L}), \mathrm{F}) = (\operatorname{Tor}_{1}^{\mathrm{R}}(\rho_{*}(\mathrm{S}_{d}), \mathrm{F}))^{(1)} = 0\) . Every right S-module E can be written in the form E = L/H for a suitable free S-module L; then we have the exact sequence

\[
0 = \operatorname{Tor} _ {1} ^ {\mathrm{R}} (\rho_ {*} (\mathrm{L}), \mathrm{F}) \rightarrow \operatorname{Tor} _ {1} ^ {\mathrm{R}} (\rho_ {*} (\mathrm{E}), \mathrm{F}) \rightarrow \rho_ {*} (\mathrm{H}) \otimes_ {\mathrm{R}} \mathrm{F} \rightarrow \rho_ {*} (\mathrm{L}) \otimes_ {\mathrm{R}} \mathrm{F}. \tag {2}
\]

But as \(\mathbf{F}_{(\mathbb{S})}\) is flat, the homomorphism \(\mathbf{H} \otimes_{\mathbb{S}} \mathbf{F}_{(\mathbb{S})} \to \mathbf{L} \otimes_{\mathbb{S}} \mathbf{F}_{(\mathbb{S})}\) is injective and is identified with the homomorphism

\[
\rho_ {*} (\mathrm{H}) \otimes_ {\mathbf {R}} \mathrm{F} \rightarrow \rho_ {*} (\mathrm{L}) \otimes_ {\mathbf {R}} \mathrm{F}.
\]

Then it follows from (2) that \(\mathrm{Tor}_1^{\mathbb{R}}(\rho_*(\mathbf{E}),\mathbf{F}) = 0\) .

Remark. Proposition 2 also follows from the existence of the exact sequence

\[
\mathrm{E} \otimes_ {\mathbf {S}} \operatorname{Tor} _ {1} ^ {\mathrm{R}} (\rho_ {*} (\mathrm{S} _ {d}), \mathrm{F}) \rightarrow \operatorname{Tor} _ {1} ^ {\mathrm{R}} (\rho_ {*} (\mathrm{E}), \mathrm{F}) \rightarrow \operatorname{Tor} _ {1} ^ {\mathrm{S}} (\mathrm{E}, \mathrm{S} _ {d} \otimes_ {\mathrm{R}} \mathrm{F}) \rightarrow 0
\]

arising from the spectral sequence of the ``associativity'' of the Tor functors.

